% Options for packages loaded elsewhere
\PassOptionsToPackage{unicode}{hyperref}
\PassOptionsToPackage{hyphens}{url}
%
\documentclass[
]{article}
\usepackage{amsmath,amssymb}
\usepackage{iftex}
\ifPDFTeX
  \usepackage[T1]{fontenc}
  \usepackage[utf8]{inputenc}
  \usepackage{textcomp} % provide euro and other symbols
\else % if luatex or xetex
  \usepackage{unicode-math} % this also loads fontspec
  \defaultfontfeatures{Scale=MatchLowercase}
  \defaultfontfeatures[\rmfamily]{Ligatures=TeX,Scale=1}
\fi
\usepackage{lmodern}
\ifPDFTeX\else
  % xetex/luatex font selection
\fi
% Use upquote if available, for straight quotes in verbatim environments
\IfFileExists{upquote.sty}{\usepackage{upquote}}{}
\IfFileExists{microtype.sty}{% use microtype if available
  \usepackage[]{microtype}
  \UseMicrotypeSet[protrusion]{basicmath} % disable protrusion for tt fonts
}{}
\makeatletter
\@ifundefined{KOMAClassName}{% if non-KOMA class
  \IfFileExists{parskip.sty}{%
    \usepackage{parskip}
  }{% else
    \setlength{\parindent}{0pt}
    \setlength{\parskip}{6pt plus 2pt minus 1pt}}
}{% if KOMA class
  \KOMAoptions{parskip=half}}
\makeatother
\usepackage{xcolor}
\usepackage[margin=1in]{geometry}
\usepackage{color}
\usepackage{fancyvrb}
\newcommand{\VerbBar}{|}
\newcommand{\VERB}{\Verb[commandchars=\\\{\}]}
\DefineVerbatimEnvironment{Highlighting}{Verbatim}{commandchars=\\\{\}}
% Add ',fontsize=\small' for more characters per line
\usepackage{framed}
\definecolor{shadecolor}{RGB}{248,248,248}
\newenvironment{Shaded}{\begin{snugshade}}{\end{snugshade}}
\newcommand{\AlertTok}[1]{\textcolor[rgb]{0.94,0.16,0.16}{#1}}
\newcommand{\AnnotationTok}[1]{\textcolor[rgb]{0.56,0.35,0.01}{\textbf{\textit{#1}}}}
\newcommand{\AttributeTok}[1]{\textcolor[rgb]{0.13,0.29,0.53}{#1}}
\newcommand{\BaseNTok}[1]{\textcolor[rgb]{0.00,0.00,0.81}{#1}}
\newcommand{\BuiltInTok}[1]{#1}
\newcommand{\CharTok}[1]{\textcolor[rgb]{0.31,0.60,0.02}{#1}}
\newcommand{\CommentTok}[1]{\textcolor[rgb]{0.56,0.35,0.01}{\textit{#1}}}
\newcommand{\CommentVarTok}[1]{\textcolor[rgb]{0.56,0.35,0.01}{\textbf{\textit{#1}}}}
\newcommand{\ConstantTok}[1]{\textcolor[rgb]{0.56,0.35,0.01}{#1}}
\newcommand{\ControlFlowTok}[1]{\textcolor[rgb]{0.13,0.29,0.53}{\textbf{#1}}}
\newcommand{\DataTypeTok}[1]{\textcolor[rgb]{0.13,0.29,0.53}{#1}}
\newcommand{\DecValTok}[1]{\textcolor[rgb]{0.00,0.00,0.81}{#1}}
\newcommand{\DocumentationTok}[1]{\textcolor[rgb]{0.56,0.35,0.01}{\textbf{\textit{#1}}}}
\newcommand{\ErrorTok}[1]{\textcolor[rgb]{0.64,0.00,0.00}{\textbf{#1}}}
\newcommand{\ExtensionTok}[1]{#1}
\newcommand{\FloatTok}[1]{\textcolor[rgb]{0.00,0.00,0.81}{#1}}
\newcommand{\FunctionTok}[1]{\textcolor[rgb]{0.13,0.29,0.53}{\textbf{#1}}}
\newcommand{\ImportTok}[1]{#1}
\newcommand{\InformationTok}[1]{\textcolor[rgb]{0.56,0.35,0.01}{\textbf{\textit{#1}}}}
\newcommand{\KeywordTok}[1]{\textcolor[rgb]{0.13,0.29,0.53}{\textbf{#1}}}
\newcommand{\NormalTok}[1]{#1}
\newcommand{\OperatorTok}[1]{\textcolor[rgb]{0.81,0.36,0.00}{\textbf{#1}}}
\newcommand{\OtherTok}[1]{\textcolor[rgb]{0.56,0.35,0.01}{#1}}
\newcommand{\PreprocessorTok}[1]{\textcolor[rgb]{0.56,0.35,0.01}{\textit{#1}}}
\newcommand{\RegionMarkerTok}[1]{#1}
\newcommand{\SpecialCharTok}[1]{\textcolor[rgb]{0.81,0.36,0.00}{\textbf{#1}}}
\newcommand{\SpecialStringTok}[1]{\textcolor[rgb]{0.31,0.60,0.02}{#1}}
\newcommand{\StringTok}[1]{\textcolor[rgb]{0.31,0.60,0.02}{#1}}
\newcommand{\VariableTok}[1]{\textcolor[rgb]{0.00,0.00,0.00}{#1}}
\newcommand{\VerbatimStringTok}[1]{\textcolor[rgb]{0.31,0.60,0.02}{#1}}
\newcommand{\WarningTok}[1]{\textcolor[rgb]{0.56,0.35,0.01}{\textbf{\textit{#1}}}}
\usepackage{graphicx}
\makeatletter
\def\maxwidth{\ifdim\Gin@nat@width>\linewidth\linewidth\else\Gin@nat@width\fi}
\def\maxheight{\ifdim\Gin@nat@height>\textheight\textheight\else\Gin@nat@height\fi}
\makeatother
% Scale images if necessary, so that they will not overflow the page
% margins by default, and it is still possible to overwrite the defaults
% using explicit options in \includegraphics[width, height, ...]{}
\setkeys{Gin}{width=\maxwidth,height=\maxheight,keepaspectratio}
% Set default figure placement to htbp
\makeatletter
\def\fps@figure{htbp}
\makeatother
\setlength{\emergencystretch}{3em} % prevent overfull lines
\providecommand{\tightlist}{%
  \setlength{\itemsep}{0pt}\setlength{\parskip}{0pt}}
\setcounter{secnumdepth}{-\maxdimen} % remove section numbering
\ifLuaTeX
  \usepackage{selnolig}  % disable illegal ligatures
\fi
\usepackage{bookmark}
\IfFileExists{xurl.sty}{\usepackage{xurl}}{} % add URL line breaks if available
\urlstyle{same}
\hypersetup{
  pdftitle={MT9: Comparing Groups},
  hidelinks,
  pdfcreator={LaTeX via pandoc}}

\title{MT9: Comparing Groups}
\author{}
\date{\vspace{-2.5em}}

\begin{document}
\maketitle

This week we will move on from the empty model to add categorical
explanatory variables that might help to reduce the error or imprecision
of statistical estimation. Let's say we have an empty model that
contained 100 student scores of narcissism. Without any other
information, our best guess of what any given student might score on
narcissism is the mean (i.e., the empty model). However, let's say we
speculate that we might be able to improve our best guess, or reduce our
estimated variance, if we knew the gender of the student because males
typically score higher on narcissism than females. This is, in essence,
what we are doing by adding a categorical explanatory variable to our
empty model -- attempting to improve the accuracy of our best guess by
adding variables that explain the error variation in the empty model.

\subsection{Learning outcomes}\label{learning-outcomes}

At the end of today's workshop you should be able to do the following:

\begin{enumerate}
\def\labelenumi{\arabic{enumi}.}
\item
  Conduct a one-parameter linear model (paired samples t-test) using the
  lm and t.test functions
\item
  Conduct a two-parameter linear model (independent samples t-test)
  using the lm and t.test model functions
\item
  Conduct a three-parameter model (between groups one-way ANOVA) using
  the lm and aov functions
\end{enumerate}

\subsection{Install required packages}\label{install-required-packages}

\begin{Shaded}
\begin{Highlighting}[]
\FunctionTok{install.packages}\NormalTok{(}\StringTok{"supernova"}\NormalTok{)}
\FunctionTok{install.packages}\NormalTok{(}\StringTok{"readr"}\NormalTok{)}
\FunctionTok{install.packages}\NormalTok{(}\StringTok{"sjlabelled"}\NormalTok{)}
\FunctionTok{install.packages}\NormalTok{(}\StringTok{"fastDummies"}\NormalTok{)}
\FunctionTok{install.packages}\NormalTok{(}\StringTok{"mosaic"}\NormalTok{)}
\FunctionTok{install.packages}\NormalTok{(}\StringTok{"rstatix"}\NormalTok{)}
\FunctionTok{install.packages}\NormalTok{(}\StringTok{"Hmisc"}\NormalTok{)}
\FunctionTok{install.packages}\NormalTok{(}\StringTok{"ggpubr"}\NormalTok{)}
\end{Highlighting}
\end{Shaded}

\subsection{Load required packages}\label{load-required-packages}

The required packages are the same as the installed packages. Write the
code needed to load the required packages in the below R chunk.

\begin{Shaded}
\begin{Highlighting}[]
\FunctionTok{library}\NormalTok{(mosaic)}
\FunctionTok{library}\NormalTok{(sjlabelled)}
\FunctionTok{library}\NormalTok{(fastDummies)}
\FunctionTok{library}\NormalTok{(supernova)}
\FunctionTok{library}\NormalTok{(readr)}
\FunctionTok{library}\NormalTok{(rstatix)}
\FunctionTok{library}\NormalTok{(Hmisc)}
\FunctionTok{library}\NormalTok{(ggpubr)}
\end{Highlighting}
\end{Shaded}

\subsection{One-parameter linear model or paired samples
t-test}\label{one-parameter-linear-model-or-paired-samples-t-test}

To practice conducting a one-parameter linear model with paired
differences (paired samples t-test), we are going to use a dataset from
Pruzek \& Helmreich (2009). This dataset presents 15 paired data
corresponding to the weights (lbs) of girls before and after 14-weeks of
psychological therapy for anorexia.

The data are from an intervention study with pre and post weights and
give a good example of a paired design. Many interventions are set up
like this in psychology, and the one-parameter linear model is how we
test for differences from pre-test to post-test.

We are going to go through testing these models with the lm() function
first, and I'll then show you the equivalence with the t.test()
function.

First, lets load this data into our R environment. Go to the MT9 folder,
and then to the workshop folder, and find the ``anorexia.csv'' file.
Click on it aand then select ``import dataset''. In the new window that
appears, click ``update'' and then when the dataframe shows, click
import. If you want, you can try running the code below and it should do
the same thing (if not put your hand up).

\begin{Shaded}
\begin{Highlighting}[]
\NormalTok{anorexia }\OtherTok{\textless{}{-}} \FunctionTok{read\_csv}\NormalTok{(}\StringTok{"C:/Users/CURRANT/Dropbox/Work/LSE/PB130/MT9/Workshop/anorexia.csv"}\NormalTok{)}
\end{Highlighting}
\end{Shaded}

\begin{verbatim}
## Rows: 15 Columns: 3
## -- Column specification --------------------------------------------------------
## Delimiter: ","
## chr (1): Subject
## dbl (2): X1, Y1
## 
## i Use `spec()` to retrieve the full column specification for this data.
## i Specify the column types or set `show_col_types = FALSE` to quiet this message.
\end{verbatim}

\begin{Shaded}
\begin{Highlighting}[]
\NormalTok{anorexia}
\end{Highlighting}
\end{Shaded}

\begin{verbatim}
## # A tibble: 15 x 3
##       X1    Y1 Subject
##    <dbl> <dbl> <chr>  
##  1  8.88 10.1  S01    
##  2 14.4  11.9  S02    
##  3  8.02  6.02 S03    
##  4  5.84  3.04 S04    
##  5  5.47  1.87 S05    
##  6 12.1  12.6  S06    
##  7 11.7   9.66 S07    
##  8 10.3   9.26 S08    
##  9  5.07  6.16 S09    
## 10  8.24 10.8  S10    
## 11 15.1  12.4  S11    
## 12 13.5  10.2  S12    
## 13 11.3  12.4  S13    
## 14  9.82  9.66 S14    
## 15  9.56  6.96 S15
\end{verbatim}

You will see that this dataframe has three columns and 15 subjects
(rows). X1 is the girls' pre-therapy weight and Y1 is the girls'
post-therapy weight. The Subject variable is the subject code for each
participant.

Before we delve into this topic lets just clarify the research question
and null hypothesis we are testing:

\emph{Research Question} - Do girls have different before and after
therapy weights?

\emph{Null Hypothesis} - The difference between girls' weight before and
after therapy will be zero.

\subsection{Descriptives}\label{descriptives}

Lets first inspect this data and take a look at the means and standard
deviations for each time point. By now, you should know how to do this.
As a reminder, we can use favstats() to inspect the descriptives.

\begin{Shaded}
\begin{Highlighting}[]
\FunctionTok{favstats}\NormalTok{(}\SpecialCharTok{\textasciitilde{}}\NormalTok{X1, }\AttributeTok{data=}\NormalTok{anorexia)}
\end{Highlighting}
\end{Shaded}

\begin{verbatim}
##    min    Q1 median    Q3   max     mean       sd  n missing
##  5.065 8.125   9.82 11.89 15.08 9.948667 3.126275 15       0
\end{verbatim}

\begin{Shaded}
\begin{Highlighting}[]
\FunctionTok{favstats}\NormalTok{(}\SpecialCharTok{\textasciitilde{}}\NormalTok{Y1, }\AttributeTok{data=}\NormalTok{anorexia)}
\end{Highlighting}
\end{Shaded}

\begin{verbatim}
##   min    Q1 median      Q3   max mean       sd  n missing
##  1.87 6.555   9.66 11.3625 12.64 8.87 3.365538 15       0
\end{verbatim}

Okay, there appears to be a small reduction in weight from pre to post
therapy evidenced by the means (not a good start). But there is quite
substantial variance in both samples, as evidenced by the SD, so the
question is whether this mean difference is statistically meaningful. Is
it large enough for us to conclude that a zero effect is unlikely?

To address this question, lets first calculate the difference scores for
each participant and take a look at the mean difference between the
pairs.

\begin{Shaded}
\begin{Highlighting}[]
\NormalTok{anorexia}\SpecialCharTok{$}\NormalTok{difference }\OtherTok{\textless{}{-}}\NormalTok{ anorexia}\SpecialCharTok{$}\NormalTok{Y1 }\SpecialCharTok{{-}}\NormalTok{ anorexia}\SpecialCharTok{$}\NormalTok{X1}
\NormalTok{anorexia}
\end{Highlighting}
\end{Shaded}

\begin{verbatim}
## # A tibble: 15 x 4
##       X1    Y1 Subject difference
##    <dbl> <dbl> <chr>        <dbl>
##  1  8.88 10.1  S01          1.25 
##  2 14.4  11.9  S02         -2.44 
##  3  8.02  6.02 S03         -1.99 
##  4  5.84  3.04 S04         -2.79 
##  5  5.47  1.87 S05         -3.6  
##  6 12.1  12.6  S06          0.58 
##  7 11.7   9.66 S07         -2.06 
##  8 10.3   9.26 S08         -1.05 
##  9  5.07  6.16 S09          1.09 
## 10  8.24 10.8  S10          2.55 
## 11 15.1  12.4  S11         -2.72 
## 12 13.5  10.2  S12         -3.31 
## 13 11.3  12.4  S13          1.08 
## 14  9.82  9.66 S14         -0.160
## 15  9.56  6.96 S15         -2.61
\end{verbatim}

Lets now use favstats again to examine the mean difference between the
paired scores.

\begin{Shaded}
\begin{Highlighting}[]
\FunctionTok{favstats}\NormalTok{(}\SpecialCharTok{\textasciitilde{}}\NormalTok{difference, }\AttributeTok{data=}\NormalTok{anorexia)}
\end{Highlighting}
\end{Shaded}

\begin{verbatim}
##   min     Q1 median   Q3  max      mean       sd  n missing
##  -3.6 -2.665  -1.99 0.83 2.55 -1.078667 1.973018 15       0
\end{verbatim}

Here we have a mean reduction of 1.08 lbs from before to after therapy.
The SD associated with that reduction is 1.97. We can also think about
this mean difference as our predicted value, like the mean was the
predicted value in the empty model. All else being equal and with no
other information, our best guess would be that the psychological
therapy would result in a weight reduction of 1.08 lbs. As we saw with
the descriptives, this is not good news for our intervention.

\subsection{Setting up the linear
model}\label{setting-up-the-linear-model}

Okay, lets set our up our linear model to examine the statistical
significance of that -1.08 lb mean difference. Like last week, we are
using the lm function to test a one-parameter model - but this time the
parameter is the difference between the pairs of observations.

\begin{Shaded}
\begin{Highlighting}[]
\NormalTok{anorexia.model }\OtherTok{\textless{}{-}} \FunctionTok{lm}\NormalTok{(Y1 }\SpecialCharTok{{-}}\NormalTok{ X1 }\SpecialCharTok{\textasciitilde{}} \ConstantTok{NULL}\NormalTok{, }\AttributeTok{data=}\NormalTok{anorexia) }\CommentTok{\# notice how I\textquotesingle{}ve set this up. The lm is taking the difference score (Y1{-} X1) and inputting it to a one{-}parameter model (NULL). I am then saving that model as an R object called anorexia.model}
\FunctionTok{summary}\NormalTok{(anorexia.model) }\CommentTok{\# the summary command just calls the summary statistics for the anorexia.model}
\end{Highlighting}
\end{Shaded}

\begin{verbatim}
## 
## Call:
## lm(formula = Y1 - X1 ~ NULL, data = anorexia)
## 
## Residuals:
##     Min      1Q  Median      3Q     Max 
## -2.5213 -1.5863 -0.9113  1.9087  3.6287 
## 
## Coefficients:
##             Estimate Std. Error t value Pr(>|t|)  
## (Intercept)  -1.0787     0.5094  -2.117   0.0526 .
## ---
## Signif. codes:  0 '***' 0.001 '**' 0.01 '*' 0.05 '.' 0.1 ' ' 1
## 
## Residual standard error: 1.973 on 14 degrees of freedom
\end{verbatim}

Notice that the model part of this formula - bo - is the same as the
mean difference we just calculated -1.08 lbs? This shows that we still
have a one-parameter model here but instead of the mean of an outcome in
the empty model, this time we have the mean of the differences in the
outcome variable for our paired t-test.

Notice also we have some other statistics that we didn't see in the
empty model. The standard error is the estimate of error in the estimate
of the model due to sampling variance. It is calculated by taking the
standard deviation of the mean difference and dividing it by the sqrt of
the sample size. Larger sample sizes and smaller mean difference
deviations result in a smaller standard error of the estimate. We
calculate this by hand using the information we have already called.

\begin{Shaded}
\begin{Highlighting}[]
\FloatTok{1.973018}\SpecialCharTok{/}\FunctionTok{sqrt}\NormalTok{(}\DecValTok{15}\NormalTok{) }\CommentTok{\# 1.97 is the SD for the mean difference and the sample size is 15}
\end{Highlighting}
\end{Shaded}

\begin{verbatim}
## [1] 0.5094311
\end{verbatim}

We also have a t-ratio. This is the ratio of the model estimate - bo -
to its standard error. The larger the estimate relative to the SE the
larger t will be. Here we have a t-ratio of -2.117. We can also
calculate this by hand using the information we have already called.

\begin{Shaded}
\begin{Highlighting}[]
\SpecialCharTok{{-}}\FloatTok{1.0787}\SpecialCharTok{/}\FloatTok{0.5094311} \CommentTok{\# ratio of the estimate (mean difference) to the standard error}
\end{Highlighting}
\end{Shaded}

\begin{verbatim}
## [1] -2.11746
\end{verbatim}

So we have a t-ratio for the mean difference between the pre and
post-therapy weights of -2.12. But is this large enough to be
statistically significant? Can we say that a zero mean difference is
unlikely given the data?

Well take a look at the p value (pr) in the lm () output and remember we
said that statisticians count .05 and lower probabilities as unlikely.
Our p value is \textgreater{} .05 and therefore we cannot reject the
null hypothesis. It seems that weight appears to fall across the time
periods in our sample but that decrease is not statistically
significant.

We might write this up in a research paper as follows:

There was no statistically significant difference between pre-therapy
and post-therapy weight in the current sample: t(14) = -2.12, p
\textgreater{} .05.

\subsection{Equivalence with the paried
t.test}\label{equivalence-with-the-paried-t.test}

Just to show you that the paired t-test is equivalent to a linear model,
lets do the same analysis using the t.test function.

\begin{Shaded}
\begin{Highlighting}[]
\FunctionTok{t.test}\NormalTok{(anorexia}\SpecialCharTok{$}\NormalTok{Y1, anorexia}\SpecialCharTok{$}\NormalTok{X1, }\AttributeTok{paired =} \ConstantTok{TRUE}\NormalTok{) }\CommentTok{\# setting up the t{-}test requires the post then pre variables entered in order rathern than as a difference score {-} R calcualtes the difference score for us in the t.test function.}
\end{Highlighting}
\end{Shaded}

\begin{verbatim}
## 
##  Paired t-test
## 
## data:  anorexia$Y1 and anorexia$X1
## t = -2.1174, df = 14, p-value = 0.05261
## alternative hypothesis: true mean difference is not equal to 0
## 95 percent confidence interval:
##  -2.17128734  0.01395401
## sample estimates:
## mean difference 
##       -1.078667
\end{verbatim}

See how the t ratio, p-value and mean difference is exactly the same as
the linear model we tested? Thats because the analyses are exactly the
same! Students are often taught paired t-test as a separate statistical
test to independent t-test, ANOVA and regression (these are tests we
will come onto here and in subsequent weeks) but, in fact, they are all
a family of tests belonging to the linear model. This equivalence helps
us understand what is happening when we run stats in the psychological
and behavioral sciences. Its all underpinned by the linear model.

\subsection{Plotting the paired sample
t-test}\label{plotting-the-paired-sample-t-test}

We might plot this difference using side-by-side boxplots like we did in
the visualization week (MT4). Here is the code to place the
distributions of the two time points side-on using ggplot.

\begin{Shaded}
\begin{Highlighting}[]
\FunctionTok{library}\NormalTok{(rstatix)}
\FunctionTok{library}\NormalTok{(ggpubr)}
\CommentTok{\# Transform into long data: }
\CommentTok{\# gather the before and after values in the same column}
\NormalTok{anorexia.long }\OtherTok{\textless{}{-}}\NormalTok{ anorexia }\SpecialCharTok{\%\textgreater{}\%}
\NormalTok{  tidyr}\SpecialCharTok{::}\FunctionTok{gather}\NormalTok{(}\AttributeTok{key =} \StringTok{"Time"}\NormalTok{, }\AttributeTok{value =} \StringTok{"weight"}\NormalTok{, X1, Y1)}

\CommentTok{\# create box{-}plots}
\NormalTok{bxp }\OtherTok{\textless{}{-}}\NormalTok{ ggpubr}\SpecialCharTok{::}\FunctionTok{ggpaired}\NormalTok{(anorexia.long, }\AttributeTok{x =} \StringTok{"Time"}\NormalTok{, }\AttributeTok{y =} \StringTok{"weight"}\NormalTok{, }
         \AttributeTok{order =} \FunctionTok{c}\NormalTok{(}\StringTok{"X1"}\NormalTok{, }\StringTok{"Y1"}\NormalTok{),}
         \AttributeTok{ylab =} \StringTok{"Weight"}\NormalTok{, }\AttributeTok{xlab =} \StringTok{"Time"}\NormalTok{)}

\CommentTok{\# Add p{-}value and significance levels}
\NormalTok{stat.test }\OtherTok{\textless{}{-}}\NormalTok{ anorexia.long  }\SpecialCharTok{\%\textgreater{}\%} 
  \FunctionTok{t\_test}\NormalTok{(weight }\SpecialCharTok{\textasciitilde{}}\NormalTok{ Time, }\AttributeTok{paired =} \ConstantTok{TRUE}\NormalTok{) }\SpecialCharTok{\%\textgreater{}\%}
  \FunctionTok{add\_significance}\NormalTok{()}
\NormalTok{stat.test }\OtherTok{\textless{}{-}}\NormalTok{ stat.test }\SpecialCharTok{\%\textgreater{}\%} \FunctionTok{add\_xy\_position}\NormalTok{(}\AttributeTok{x =} \StringTok{"Time"}\NormalTok{)}
\NormalTok{bxp }\SpecialCharTok{+} \FunctionTok{stat\_pvalue\_manual}\NormalTok{(stat.test, }\AttributeTok{tip.length =} \DecValTok{0}\NormalTok{) }\SpecialCharTok{+}
  \FunctionTok{labs}\NormalTok{(}\AttributeTok{subtitle =} \FunctionTok{get\_test\_label}\NormalTok{(stat.test, }\AttributeTok{detailed=} \ConstantTok{TRUE}\NormalTok{))}
\end{Highlighting}
\end{Shaded}

\includegraphics{Notebook_empty_files/figure-latex/unnamed-chunk-10-1.pdf}

\subsection{How it's written}\label{how-its-written}

We could report the result as follows:

The mean pre-intervention weight was 9.95 lbs (SD = 3.13), whereas the
mean post-intervention weight was 8.87 lbs(SD = 3.37). A paired-sample
t-test showed that the difference was not statistically significant,
mean diff = -1.08 lbs, t(14) = -2.12, p \textgreater{} 0.05.

\subsection{Two-parameter model or independent
t-test}\label{two-parameter-model-or-independent-t-test}

Okay lets now move to a different research topic and question to examine
differences between independent groups. We will start with a
two-parameter model (two groups) and move to a three parameter model. As
we saw in the lecture, the two-parameter model is often called an
independent t-test and the three-parameter model is often called ANOVA.
Like the paired t-test, we will see that both of these tests are linear
models that can be analyzed using the lm() in R.

To start, lets ask different research question using some data I have
collected on levels of perfectionism in American, Canadian, and British
college students. This data is from Curran \& Hill (2019). Many
observational studies in psychology are set up like this, and the
two-parameter linear model is how we test for differences in scores
across the countries.

Lets start by collapsing the USA and Canada into one region called North
America and look to see if North American students differ in their
levels of perfectionism to those in the UK. Later, we will extend this
model to look at differences across the three counties.

First, lets load this data into our R environment. Go to the MT9 folder,
and then to the workshop folder, and find the ``perfectionism.csv''
file. Click on it and then select ``import dataset''. In the new window
that appears, click ``update'' and then when the dataframe shows, click
import. If you want, you can run the code below and it will do the same
thing.

\begin{Shaded}
\begin{Highlighting}[]
\FunctionTok{library}\NormalTok{(readr)}
\NormalTok{perfectionism }\OtherTok{\textless{}{-}} \FunctionTok{read\_csv}\NormalTok{(}\StringTok{"C:/Users/CURRANT/Dropbox/Work/LSE/PB130/MT9/Workshop/perfectionism.csv"}\NormalTok{)}
\end{Highlighting}
\end{Shaded}

\begin{verbatim}
## Rows: 139 Columns: 4
## -- Column specification --------------------------------------------------------
## Delimiter: ","
## chr (3): Study, Country, Region
## dbl (1): SOP
## 
## i Use `spec()` to retrieve the full column specification for this data.
## i Specify the column types or set `show_col_types = FALSE` to quiet this message.
\end{verbatim}

\begin{Shaded}
\begin{Highlighting}[]
\NormalTok{perfectionism}
\end{Highlighting}
\end{Shaded}

\begin{verbatim}
## # A tibble: 139 x 4
##    Study            SOP Country Region       
##    <chr>          <dbl> <chr>   <chr>        
##  1 Hewitt & Flett  4.39 CAN     NORTH_AMERICA
##  2 Hewitt & Flett  4.53 CAN     NORTH_AMERICA
##  3 Hewitt & Flett  4.45 CAN     NORTH_AMERICA
##  4 Westra TH       4.64 CAN     NORTH_AMERICA
##  5 Flett           4.57 CAN     NORTH_AMERICA
##  6 Flett           4.56 CAN     NORTH_AMERICA
##  7 Flett           4.53 CAN     NORTH_AMERICA
##  8 Flett           4.67 CAN     NORTH_AMERICA
##  9 Hart et al      4.92 UK      UK           
## 10 Bottos & Dewey  4.81 CAN     NORTH_AMERICA
## # i 129 more rows
\end{verbatim}

You will see that this dataframe has four columns and 139 data points
(rows). Study is the study that the mean level of perfectionism was
taken from, SOP is the mean level of self-oriented perfectionism
reported in the study, Country is the varaible containing the country
that the data was collected in (i.e., Canada, the USA, and the UK), and
Region (i.e., North America and the UK) is the World region that the
data was collected from.

Before we delve into this topic lets just clarify the research question
for our two-parameter model and null hypothesis we are testing:

\emph{Research Question} - Do levels of perfectionism differ between
college students from North America and the UK?

\emph{Null Hypothesis} - The difference between college students' levels
of perfectionism in North America and the UK will be zero.

\subsection{Descriptives}\label{descriptives-1}

Lets first inspect this data and take a look at the means and standard
devations for each region. By now, you should know how to do this. As a
reminder, we can use favstats() to inpect the descriptives.

\begin{Shaded}
\begin{Highlighting}[]
\FunctionTok{favstats}\NormalTok{(}\SpecialCharTok{\textasciitilde{}}\NormalTok{SOP, Region, }\AttributeTok{data=}\NormalTok{perfectionism) }\CommentTok{\# the second line of fav stats is our factor and it splits the SOP scores into a mean for North America and a mean for the UK.}
\end{Highlighting}
\end{Shaded}

\begin{verbatim}
##          Region      min       Q1   median     Q3  max     mean        sd   n
## 1 NORTH_AMERICA 4.140000 4.480000 4.595333 4.7190 5.22 4.607002 0.1861641 115
## 2            UK 3.576667 4.335333 4.584333 4.7445 5.21 4.513348 0.3505012  24
##   missing
## 1       0
## 2       0
\end{verbatim}

Okay, there appears to be a small difference here and we can calculate
this difference by hand. We just take the UK mean from the North America
mean to ascertain the increment needed from the UK to get the North
America mean.

\begin{Shaded}
\begin{Highlighting}[]
\FloatTok{4.607002} \SpecialCharTok{{-}} \FloatTok{4.513348}
\end{Highlighting}
\end{Shaded}

\begin{verbatim}
## [1] 0.093654
\end{verbatim}

North American students appear to have a higher scores on perfectionism
than UK students. The difference is 0.093654 points on a 1-7 Likert
scale, which doesnt seem a great deal but in the context of the equally
small SD, could be significant. The question therefore is whether this
mean difference is statistically meaningful or large enough for us to
consider the possibility of a zero difference in the population
unlikely.

\subsection{Setting up the linear
model}\label{setting-up-the-linear-model-1}

Now lets set our up our linear model to examine the statitical
significance of that 0.09 mean difference. This time though we are using
the lm() function to test a two-parameter linear model becuase we now
have two means to estimate.

Before we do so, I want to show you how to dummy code the variables
because the meaning of the intercept changes in line with how the
explanatory variable is coded. The intercept, remember, is the value of
the outcome (perfectionism) when the explanatory variable equals zero.
And so the meaning of zero becomes very important. In this example, I
want to know what the increment is from the UK mean to the North America
mean, so I'm going to dummy code my explanatory variable as UK = 0 and
North America = 1. This way, the interpretation of b0 (intercept) is the
UK mean, and the interpretation of b1 is the increment needed for the
North America mean or - in other words - the mean difference. To do so,
I'm just going use the ifelse() function to assign codes.

\begin{Shaded}
\begin{Highlighting}[]
\NormalTok{perfectionism}\SpecialCharTok{$}\NormalTok{dummy }\OtherTok{\textless{}{-}} \FunctionTok{ifelse}\NormalTok{(perfectionism}\SpecialCharTok{$}\NormalTok{Region }\SpecialCharTok{==} \FunctionTok{c}\NormalTok{(}\StringTok{"NORTH\_AMERICA"}\NormalTok{), }\DecValTok{1}\NormalTok{, }\DecValTok{0}\NormalTok{) }\CommentTok{\# this code creates new variable "dummy" and inserts either 1 or 0 depending on whether North America is TRUE in the Region variable}
\NormalTok{perfectionism}
\end{Highlighting}
\end{Shaded}

\begin{verbatim}
## # A tibble: 139 x 5
##    Study            SOP Country Region        dummy
##    <chr>          <dbl> <chr>   <chr>         <dbl>
##  1 Hewitt & Flett  4.39 CAN     NORTH_AMERICA     1
##  2 Hewitt & Flett  4.53 CAN     NORTH_AMERICA     1
##  3 Hewitt & Flett  4.45 CAN     NORTH_AMERICA     1
##  4 Westra TH       4.64 CAN     NORTH_AMERICA     1
##  5 Flett           4.57 CAN     NORTH_AMERICA     1
##  6 Flett           4.56 CAN     NORTH_AMERICA     1
##  7 Flett           4.53 CAN     NORTH_AMERICA     1
##  8 Flett           4.67 CAN     NORTH_AMERICA     1
##  9 Hart et al      4.92 UK      UK                0
## 10 Bottos & Dewey  4.81 CAN     NORTH_AMERICA     1
## # i 129 more rows
\end{verbatim}

Now a new variable is created that assigns UK a code 0 and North America
a code 1. In reality, the lm() function in R will do this for you if you
feed it a categorical variable, but I wanted to show you how this is
working.

Okay lets now build our two-parameter model.

\begin{Shaded}
\begin{Highlighting}[]
\NormalTok{perfectionism.model }\OtherTok{\textless{}{-}} \FunctionTok{lm}\NormalTok{(SOP }\SpecialCharTok{\textasciitilde{}} \DecValTok{1} \SpecialCharTok{+}\NormalTok{ dummy, }\AttributeTok{data=}\NormalTok{perfectionism) }\CommentTok{\# notice how I\textquotesingle{}ve set this up. Unlike the empty or one{-}parameter model we now have a variable after the tilda (\textasciitilde{}). This is no longer an empty model with just the mean or intercept but it now contains an intercept (1) and an explainatory variable {-} namely, our dummy coded region. Thus y \textasciitilde{} 1 + x is the R{-}way of writing y = b + mx (the stright line equation!).}

\FunctionTok{summary}\NormalTok{(perfectionism.model) }\CommentTok{\# the summary command just calls the summary statistics for the perfectionism.model}
\end{Highlighting}
\end{Shaded}

\begin{verbatim}
## 
## Call:
## lm(formula = SOP ~ 1 + dummy, data = perfectionism)
## 
## Residuals:
##      Min       1Q   Median       3Q      Max 
## -0.93668 -0.13667  0.00833  0.13300  0.69665 
## 
## Coefficients:
##             Estimate Std. Error t value Pr(>|t|)    
## (Intercept)  4.51335    0.04540  99.417   <2e-16 ***
## dummy        0.09365    0.04991   1.876   0.0627 .  
## ---
## Signif. codes:  0 '***' 0.001 '**' 0.01 '*' 0.05 '.' 0.1 ' ' 1
## 
## Residual standard error: 0.2224 on 137 degrees of freedom
## Multiple R-squared:  0.02506,    Adjusted R-squared:  0.01794 
## F-statistic: 3.521 on 1 and 137 DF,  p-value: 0.06272
\end{verbatim}

Notice that the intercept part of our model - bo - is the same as the
mean perfectionism score for the UK - 4.51? Notice also that b1 is the
mean difference or the increment needed to get to the North American
mean (i.e., 0.09 + 4.51 = 4.60)? This shows that we now have two
parameters, and a mean difference that is exactly as we calculated by
hand.

Notice also we have the same other statistics as in our previous model.
The standard error and the t-ratio. From what you now know about these
statistics, you should be able to work out whether the difference of .09
is statistically meaningful. We have a t of 1.88 and a p value of .06 -
we cannot reject the null hypothesis.

\subsection{Partitioning variance}\label{partitioning-variance}

Another, perhaps more informative, way to think about this analysis is
that we are adding an explanatory variable to the empty model. By doing
this, we are attempting to reducing the empty model error. Or, in other
words, to explain or reduce the empty model variance. This is the key to
statistics; entering information to models that helps us make better
predictions. Our reasoning here is that we can better predict
perfectionism scores if we know which region a person comes from.

We've established that we cannot be certain that the mean difference in
the population is not zero (i.e., p \textgreater{} .05), but how much
variance in perfectionism is explained by region? We can answer this
question using the supernova() function in R, which breaks down the
variance in perfectionism due to the model (ie., Region) and error
(i.e., variance left over once we've subtracted out the model)

\begin{Shaded}
\begin{Highlighting}[]
\FunctionTok{library}\NormalTok{(supernova)}
\FunctionTok{supernova}\NormalTok{(perfectionism.model)}
\end{Highlighting}
\end{Shaded}

\begin{verbatim}
##  Analysis of Variance Table (Type III SS)
##  Model: SOP ~ 1 + dummy
## 
##                             SS  df    MS     F   PRE     p
##  ----- --------------- | ----- --- ----- ----- ----- -----
##  Model (error reduced) | 0.174   1 0.174 3.521 .0251 .0627
##  Error (from model)    | 6.776 137 0.049                  
##  ----- --------------- | ----- --- ----- ----- ----- -----
##  Total (empty model)   | 6.951 138 0.050
\end{verbatim}

Taking a look at the supernova table tells us that the total variance -
or sum of squares - in perfectionism scores is 6.95. This is variance of
perfectionism in the empty model (i.e., the mean as the model). Entering
our predictor in the model reduced the sum of squares in the empty
perfectionism model by 0.17. We have, then, a sum of squares 6.78 left
after subtracting out the model that includes our predictor (i.e.,
6.95-0.17).

When we look at the ratio of total variance (in the empty model) to
error variance (in the one-parameter model) we are left with the
proportion of empty model variance explained by the model, or PRE. We
can see in this example that the PRE is .03. So only about 3\% of
perfectionism variance is explained by Region. Not a great deal.

Everytime we add a predictor to the model, we spend a degree of freedom
and the RPE does not take this into consideration (because it is
calculated using the sum of squares). However, the F ratio does take
degrees of freedom into consideration because it is the ratio of model
variance (SS/df) to error variance (SS/df). You can think about it like
the ratio variance explained in the model to variance unexplained in the
model. Because variance is calculated using the degrees of freedom, the
F-ratio tells us whether the degrees of freedom that we spent in order
to make our model more complicated (i.e., adding a parameter) were
``worth it''.

In this case, we see an F ratio of 3.52 and a p value of .06. In line
with what you understand about p values, this is not statistically
significant. Therefore, the interpretation is that adding region to the
empty model was not worth it. The variance explained was too small to
justify the additional parameter. Notice that this chimes with the
interpretation of the t ratio - that there is no significant difference
between the groups.

\subsection{Equivalence with the independent
t.test}\label{equivalence-with-the-independent-t.test}

Just to show you that the independent t-test is equivalent to a linear
model, lets do the same analysis using the t.test function.

\begin{Shaded}
\begin{Highlighting}[]
\FunctionTok{t.test}\NormalTok{(perfectionism}\SpecialCharTok{$}\NormalTok{SOP }\SpecialCharTok{\textasciitilde{}}\NormalTok{ perfectionism}\SpecialCharTok{$}\NormalTok{dummy, }\AttributeTok{var.equal=}\ConstantTok{TRUE}\NormalTok{) }\CommentTok{\# setting up the t{-}test requires the outcome variance and the explanatory variance separated by tilda (\textasciitilde{}) {-} R calcualtes the difference score for us in the t.test function.}
\end{Highlighting}
\end{Shaded}

\begin{verbatim}
## 
##  Two Sample t-test
## 
## data:  perfectionism$SOP by perfectionism$dummy
## t = -1.8764, df = 137, p-value = 0.06272
## alternative hypothesis: true difference in means between group 0 and group 1 is not equal to 0
## 95 percent confidence interval:
##  -0.192349245  0.005041167
## sample estimates:
## mean in group 0 mean in group 1 
##        4.513348        4.607002
\end{verbatim}

See how the t ratio, p-value and mean difference is exactly the same as
the linear model we tested? Thats because the analyses are exactly the
same! Students are often taught independent t-test as a separate
statistical test to paired t-test, ANOVA, and regression but, in fact,
they are all a family of tests belonging to the linear model! Its all
underpinned by the linear model.

\subsection{Effect size}\label{effect-size}

The effect size is calculated by dividing the mean difference between
the groups by the pooled standard deviation.

Cohen's d formula:

d = (mean1 - mean2)/pooled.sd, where:

pooled.sd is the common standard deviation of the two groups.

pooled.sd = sqrt({[}var1\emph{(n1-1) + var2}(n2-1){]}/{[}n1 + n2 -2{]});

var1 and var2 are the variances (squared standard deviation) of group1
and 2, respectively.

n1 and n2 are the sample counts for group 1 and 2, respectively. mean1
and mean2 are the means of each group, respectively.

We can do this by hand. But we can also calculate this directly (and
more quickly) using the cohens\_d function, which is part of the package
rstatix:

\begin{Shaded}
\begin{Highlighting}[]
\NormalTok{perfectionism }\SpecialCharTok{\%\textgreater{}\%}
\FunctionTok{cohens\_d}\NormalTok{(SOP }\SpecialCharTok{\textasciitilde{}}\NormalTok{ Region)}
\end{Highlighting}
\end{Shaded}

\begin{verbatim}
## # A tibble: 1 x 7
##   .y.   group1        group2 effsize    n1    n2 magnitude
## * <chr> <chr>         <chr>    <dbl> <int> <int> <ord>    
## 1 SOP   NORTH_AMERICA UK       0.334   115    24 small
\end{verbatim}

According to Cohen (1988, 1992), the effect size is low if the value of
d varies around 0.2, medium if d varies around 0.5, and large if d
varies more than 0.8. An d of .33 is a small effect.

\subsection{Plotting the independent
t-test}\label{plotting-the-independent-t-test}

We might plot this difference using side-by-side boxplots like we did in
the visualization week (MT4). Here is the code to place the
distributions of the two categories side-on using ggplot.

\begin{Shaded}
\begin{Highlighting}[]
\NormalTok{perfectionism }\SpecialCharTok{\%\textgreater{}\%}
\FunctionTok{ggplot}\NormalTok{(}\FunctionTok{aes}\NormalTok{(}\AttributeTok{x=}\NormalTok{ Region, }\AttributeTok{y =}\NormalTok{ SOP)) }\SpecialCharTok{+} \CommentTok{\# remember to use the category variable and not the dummy variable as x here}
  \FunctionTok{geom\_boxplot}\NormalTok{(}\AttributeTok{width =} \FloatTok{0.3}\NormalTok{, }\AttributeTok{fill =} \StringTok{"white"}\NormalTok{) }\SpecialCharTok{+}
  \FunctionTok{geom\_jitter}\NormalTok{(}\FunctionTok{aes}\NormalTok{(}\AttributeTok{color =} \FunctionTok{as.factor}\NormalTok{(Region), }\AttributeTok{shape =} \FunctionTok{as.factor}\NormalTok{(Region)), }\AttributeTok{width =} \FloatTok{0.1}\NormalTok{, }\AttributeTok{size =} \DecValTok{1}\NormalTok{) }\SpecialCharTok{+} 
  \FunctionTok{xlab}\NormalTok{(}\ConstantTok{NULL}\NormalTok{) }\SpecialCharTok{+}
  \FunctionTok{ggtitle}\NormalTok{(}\StringTok{"Box{-}Plots for Perfectionism in UK and North America"}\NormalTok{) }\SpecialCharTok{+}
  \FunctionTok{theme\_minimal}\NormalTok{(}\AttributeTok{base\_size =} \DecValTok{8}\NormalTok{) }\SpecialCharTok{+}
  \FunctionTok{theme}\NormalTok{(}\AttributeTok{legend.title=}\FunctionTok{element\_blank}\NormalTok{())}
\end{Highlighting}
\end{Shaded}

\includegraphics{Notebook_empty_files/figure-latex/unnamed-chunk-19-1.pdf}

\subsection{How its written}\label{how-its-written-1}

We could report the result as follow:

The mean perfectionism in UK group was 4.51 (SD = .19), whereas the mean
in North America group was 4.61(SD = .35). A two-samples t-test showed
that the difference was not statistically significant, t(137) = -1.8764,
p \textgreater{} 0.05, d = .33; where, t(137) is shorthand notation for
a t-statistic that has 137 degrees of freedom.

\subsection{Three-parameter model or one-way
ANOVA}\label{three-parameter-model-or-one-way-anova}

Lets stay with the perfectionism data but this time we are going to test
for difference between USA, Canada, and the UK rather than the UK and
North America. Such an analysis, whereby three means are compared is
often called a between groups one-way ANOVA in the literature, but we
will see it is really just another linear model - this time with three
parameters rather than two.

Before we start, lets just clarify the research question for our
three-parameter linear model and null hypothesis we are testing:

\emph{Research Question} - Do levels of perfectionism differ between
college students from the USA, Canada, and the UK?

\emph{Null Hypothesis} - The difference between college students' levels
of perfectionism in the USA, Canada, and UK will be zero.

\subsection{Descriptives}\label{descriptives-2}

Lets first inspect this data and take a look at the means and standard
devations for each country. By now, you should know how to do this. As a
reminder, we can use favstats() to inpect the descriptives.

\begin{Shaded}
\begin{Highlighting}[]
\FunctionTok{favstats}\NormalTok{(}\SpecialCharTok{\textasciitilde{}}\NormalTok{SOP, Country, }\AttributeTok{data=}\NormalTok{perfectionism) }\CommentTok{\# the second line of fav stats is our factor and it splits the SOP scores into a mean for USA, a mean for Canada, and a mean for the UK.}
\end{Highlighting}
\end{Shaded}

\begin{verbatim}
##   Country      min       Q1   median     Q3      max     mean        sd  n
## 1     CAN 4.140000 4.446000 4.569333 4.6940 5.220000 4.576998 0.1961644 73
## 2      UK 3.576667 4.335333 4.584333 4.7445 5.210000 4.513348 0.3505012 24
## 3     USA 4.368000 4.578667 4.685216 4.7450 4.997333 4.659151 0.1562029 42
##   missing
## 1       0
## 2       0
## 3       0
\end{verbatim}

Okay, there appears to be small differences here and we can calculate
these differences by hand.

\begin{Shaded}
\begin{Highlighting}[]
\FloatTok{4.576998{-}4.513348} \CommentTok{\# CAN{-}UK diff}
\end{Highlighting}
\end{Shaded}

\begin{verbatim}
## [1] 0.06365
\end{verbatim}

\begin{Shaded}
\begin{Highlighting}[]
\FloatTok{4.576998{-}4.659151} \CommentTok{\# CAN{-}USA diff}
\end{Highlighting}
\end{Shaded}

\begin{verbatim}
## [1] -0.082153
\end{verbatim}

\begin{Shaded}
\begin{Highlighting}[]
\FloatTok{4.513348{-}4.659151} \CommentTok{\# UK{-}USA diff}
\end{Highlighting}
\end{Shaded}

\begin{verbatim}
## [1] -0.145803
\end{verbatim}

We can see that Canadians have higher levels of perfectionism than Brits
but lower levels than Americans. Americans have higher levels of
perfectionism than Brits. Again, the differences are small, but in the
context of the small SD these may be significant. The question therefore
is whether these mean differences are statistically meaningful.

\subsection{Setting up the linear
model}\label{setting-up-the-linear-model-2}

Now lets set our up our linear model to examine the statistical
significance of those mean differences. This time though we are using
the lm() function to test a three-parameter model because we now have
three means to estimate.

Before we do so, like the two-parameter model, I want to show you how to
dummy code the variables in a three-group model because the meaning of
zero is very important. In this case, we need to create two dummy
variables with UK as the reference group. I have chosen UK as the
reference group because we know it has the lowest mean (but you can set
it up whichever way you want). This necessitates the creation of a USA
vs others contrast (USA = 1, others = 0), and a Canada vs others
contrast (Canada = 1, others = 0). Setting it up this way makes in the
interpretation of the intercept the UK mean (because UK is always coded
0), b1 the increment needed to find the USA mean from the UK mean, and
b2 the increment needed to find the Canada mean from the UK mean. The
difference between b1 and b2 gives the mean difference between the USA
and Canada.

We could code the variables by hand but there is a useful function in R
that does this for us. Its called fastdummies and uses a dummy\_cols()
function to dummy code variables given a categorical input. So lets go
ahead and create those dummy variables.

\begin{Shaded}
\begin{Highlighting}[]
\NormalTok{perfectionism }\OtherTok{\textless{}{-}} \FunctionTok{dummy\_cols}\NormalTok{(perfectionism, }\AttributeTok{select\_columns =} \StringTok{"Country"}\NormalTok{) }\CommentTok{\# into the perfectionism dataframe create dummy columns using the "Country" variable in the perfectionism data frame}
\NormalTok{perfectionism}
\end{Highlighting}
\end{Shaded}

\begin{verbatim}
## # A tibble: 139 x 8
##    Study            SOP Country Region  dummy Country_CAN Country_UK Country_USA
##    <chr>          <dbl> <chr>   <chr>   <dbl>       <int>      <int>       <int>
##  1 Hewitt & Flett  4.39 CAN     NORTH_~     1           1          0           0
##  2 Hewitt & Flett  4.53 CAN     NORTH_~     1           1          0           0
##  3 Hewitt & Flett  4.45 CAN     NORTH_~     1           1          0           0
##  4 Westra TH       4.64 CAN     NORTH_~     1           1          0           0
##  5 Flett           4.57 CAN     NORTH_~     1           1          0           0
##  6 Flett           4.56 CAN     NORTH_~     1           1          0           0
##  7 Flett           4.53 CAN     NORTH_~     1           1          0           0
##  8 Flett           4.67 CAN     NORTH_~     1           1          0           0
##  9 Hart et al      4.92 UK      UK          0           0          1           0
## 10 Bottos & Dewey  4.81 CAN     NORTH_~     1           1          0           0
## # i 129 more rows
\end{verbatim}

You will see that this function returns three new dummy variables -
Country\_CAN, Country\_USA, and Country\_UK. This is quite cool because
it is giving us the option of which group to use as the reference. I
have said we should use the UK as the reference because it has the
lowest mean. So the two variables we will carry forward for our analysis
are: Country\_CAN and Country\_USA.

Okay lets now build our three-parameter model.

\begin{Shaded}
\begin{Highlighting}[]
\NormalTok{perfectionism.model2 }\OtherTok{\textless{}{-}} \FunctionTok{lm}\NormalTok{(SOP }\SpecialCharTok{\textasciitilde{}} \DecValTok{1} \SpecialCharTok{+}\NormalTok{ Country\_USA }\SpecialCharTok{+}\NormalTok{ Country\_CAN, }\AttributeTok{data=}\NormalTok{perfectionism) }\CommentTok{\# notice how I\textquotesingle{}ve set this up. Unlike the empty or two{-}parameter models we now have two variables after the tilda (\textasciitilde{}). This is no longer an empty model with just the mean or intercept and one additional parameter, but it now contains an intercept (1) and a two explainatory variables {-} namely, our dummy coded coutries USA and Canada.}

\FunctionTok{summary}\NormalTok{(perfectionism.model2) }\CommentTok{\# the summary command just calls the summary statistics for the perfectionism.model}
\end{Highlighting}
\end{Shaded}

\begin{verbatim}
## 
## Call:
## lm(formula = SOP ~ 1 + Country_USA + Country_CAN, data = perfectionism)
## 
## Residuals:
##      Min       1Q   Median       3Q      Max 
## -0.93668 -0.13141  0.01367  0.12900  0.69665 
## 
## Coefficients:
##             Estimate Std. Error t value Pr(>|t|)    
## (Intercept)  4.51335    0.04496 100.396   <2e-16 ***
## Country_USA  0.14580    0.05635   2.587   0.0107 *  
## Country_CAN  0.06365    0.05182   1.228   0.2215    
## ---
## Signif. codes:  0 '***' 0.001 '**' 0.01 '*' 0.05 '.' 0.1 ' ' 1
## 
## Residual standard error: 0.2202 on 136 degrees of freedom
## Multiple R-squared:  0.05094,    Adjusted R-squared:  0.03699 
## F-statistic:  3.65 on 2 and 136 DF,  p-value: 0.02856
\end{verbatim}

Notice that the intercept part of our model - bo - is the same as the
mean perfectionism score for the UK - 4.51? Notice also that b1 is the
mean difference or the increment needed to get to the USA mean (i.e.,
0.15 + 4.51 = 4.66) and b2 is the mean difference or increment needed to
get to the Canada mean (i.e., 0.06 + 4.51 = 4.57)? The difference
between b1 and b2 (.08) is the mean difference between USA and Canada.
This shows that we now have three parameters, and mean differences that
are exactly as we calculated by hand.

Notice also we have the same other statistics as in our previous model.
The standard error and the t-ratio. From what you now know about these
statistics, you should be able to work out whether the difference of .15
between the UK and USA is statistically significant. We have a t of and
a p value of .01 - we can reject the null hypothesis. For the UK vs
Canada difference of .06 we have a t ratio of 1.228 and a p value of
0.22 - we cannot reject he null hypothesis.

\subsection{Plotting the one-way
ANOVA}\label{plotting-the-one-way-anova}

We might plot this difference using side-by-side boxplots like we did
above. Here is the code to place the distributions of the three
categories side-on using ggplot.

\begin{Shaded}
\begin{Highlighting}[]
\NormalTok{perfectionism }\SpecialCharTok{\%\textgreater{}\%}
\FunctionTok{ggplot}\NormalTok{(}\FunctionTok{aes}\NormalTok{(}\AttributeTok{x=}\NormalTok{ Country, }\AttributeTok{y =}\NormalTok{ SOP)) }\SpecialCharTok{+} \CommentTok{\# remember to use the category variable and not the dummy variable as x here}
  \FunctionTok{geom\_boxplot}\NormalTok{(}\AttributeTok{width =} \FloatTok{0.3}\NormalTok{, }\AttributeTok{fill =} \StringTok{"white"}\NormalTok{) }\SpecialCharTok{+}
  \FunctionTok{geom\_jitter}\NormalTok{(}\FunctionTok{aes}\NormalTok{(}\AttributeTok{color =} \FunctionTok{as.factor}\NormalTok{(Country), }\AttributeTok{shape =} \FunctionTok{as.factor}\NormalTok{(Country)), }\AttributeTok{width =} \FloatTok{0.1}\NormalTok{, }\AttributeTok{size =} \DecValTok{1}\NormalTok{) }\SpecialCharTok{+} 
  \FunctionTok{xlab}\NormalTok{(}\ConstantTok{NULL}\NormalTok{) }\SpecialCharTok{+}
  \FunctionTok{ggtitle}\NormalTok{(}\StringTok{"Box{-}Plots for Perfectionism in UK, Canada, and USA"}\NormalTok{) }\SpecialCharTok{+}
  \FunctionTok{theme\_minimal}\NormalTok{(}\AttributeTok{base\_size =} \DecValTok{8}\NormalTok{) }\SpecialCharTok{+}
  \FunctionTok{theme}\NormalTok{(}\AttributeTok{legend.title=}\FunctionTok{element\_blank}\NormalTok{())}
\end{Highlighting}
\end{Shaded}

\includegraphics{Notebook_empty_files/figure-latex/unnamed-chunk-24-1.pdf}

\subsection{Partitioning variance}\label{partitioning-variance-1}

As we saw with the two-parameter model, another way to think about this
analysis is that we are adding two explainatory variables to the empty
model. By doing this, we are attempting to reduce the empty model error.
Or, in other words, to explain or reduce the empty model variance. That
is the key to statistics; entering information to models that helps us
make better predictions. Our reasoning here is that we can better
predict perfectionism scores if we know which country a person comes
from.

We've established that the mean difference in the population is not zero
for USA vs UK but not Canada vs UK, but how much variance in
perfectionism is explained by country? And is this meaningful? Again, we
can answer this question using the supernova() function in R, which
breaks down the variance in perfectionism due to the model (ie.,
country) and error (i.e., variance left over once we've subtracted out
the model).

\begin{Shaded}
\begin{Highlighting}[]
\FunctionTok{supernova}\NormalTok{(perfectionism.model2)}
\end{Highlighting}
\end{Shaded}

\begin{verbatim}
##  Analysis of Variance Table (Type III SS)
##  Model: SOP ~ 1 + Country_USA + Country_CAN
## 
##                                   SS  df    MS     F   PRE     p
##  ----------- --------------- | ----- --- ----- ----- ----- -----
##        Model (error reduced) | 0.354   2 0.177 3.650 .0509 .0286
##  Country_USA                 | 0.325   1 0.325 6.694 .0469 .0107
##  Country_CAN                 | 0.073   1 0.073 1.509 .0110 .2215
##        Error (from model)    | 6.597 136 0.049                  
##  ----------- --------------- | ----- --- ----- ----- ----- -----
##        Total (empty model)   | 6.951 138 0.050
\end{verbatim}

Taking a look at the supernova table tells us that the total variance -
or sum of squares - in perfectionism scores is 6.95. This is variance of
perfectionism in the empty model (i.e., the mean as the model). Entering
our two dummy predictors in the model reduced the sum of squares in the
empty perfectionism model by 0.354. We have, then, a sum of squares 6.60
left after subtracting out the model that includes our predictor (i.e.,
6.95-0.35). Compare this with the meager 0.17 sum of squares explained
by Region - splitting region in two and looking at perfectionism scores
across the USA, Canada, and UK has explained far more variance. Why?
Well, because Americans have much higher perfectionism scores than Brits
and Canadians! In other words, the USA explains error variance in the
empty perfectionism model.

When we look at the ratio of total variance (in the empty model) to
error variance (in the one-parameter model) we are left with the
proportion of empty model variance explained by the model, or PRE. We
can see in this example that the PRE is .05. So about 5\% of
perfectionism variance is explained by Country. Again, compare this to
the 3\% explained by Region in the two-parameter model.

Remember, though, that adding an additional parameter will always result
in some variance explained by chance, so the question becomes - is this
increase in explained variance enough given the additional model
complexity?

Everytime we add a predictor to the model, we spend a degree of freedom
and the RPE does not take this into consideration. However, the F ratio
does because it is the ratio of model variance to error variance.
Because variance is calculated using the degrees of freedom, the F-ratio
tells us whether the degrees of freedom that we spent in order to make
our model more complicated (i.e., adding a parameter) as ``worth it''.

In this case, we see an F ratio associated with the model of 3.65 and a
p value of .03. In line with what you understand about p values, this is
a statistically significant result! Therefore, the interpretation is
that adding our two Country dummy variables to the empty model was worth
it. The variance explained is large enough to justify the additional
parameter. Notice that this chimes with the interpretation of the t
ratio - that there was a significant difference between the groups.

\subsection{Equivalence to one-way
ANOVA}\label{equivalence-to-one-way-anova}

Just to show you that the between groups one-way between-person
(independent) ANOVA is equivalent to a linear model, lets do the same
analysis using the anova function.

\begin{Shaded}
\begin{Highlighting}[]
\NormalTok{perfectionism.ANOVA }\OtherTok{\textless{}{-}} \FunctionTok{aov}\NormalTok{(SOP }\SpecialCharTok{\textasciitilde{}}\NormalTok{ Country, }\AttributeTok{data =}\NormalTok{ perfectionism) }\CommentTok{\# setting up the ANOVA requires the outcome variable and the explanatory variable separated by tilde (\textasciitilde{}).}
\FunctionTok{summary}\NormalTok{(perfectionism.ANOVA) }
\end{Highlighting}
\end{Shaded}

\begin{verbatim}
##              Df Sum Sq Mean Sq F value Pr(>F)  
## Country       2  0.354  0.1771    3.65 0.0286 *
## Residuals   136  6.597  0.0485                 
## ---
## Signif. codes:  0 '***' 0.001 '**' 0.01 '*' 0.05 '.' 0.1 ' ' 1
\end{verbatim}

See how the F ratio and mean differences are the same as the linear
model we tested? Thats because the analyses are exactly the same!
Students are often taught one-way ANOVA as a separate statistical test
to paired t-test, independent t-test, and regression (these are tests we
will come onto here and in subsequent weeks) but, in fact, they are all
a family of tests belonging to the linear model!

As it turns out, its better to run one-way ANOVA analyses with aov()
because the number of comparisons being made is more than 2 and
therefore dummy variables are needed. As the number of comparisons
increase from 3, to 4, to 5, and so on the number of dummy variables
also increases. ANOVA solves this by running whats called the omnibus F
test, which is one test to capture the variance explained in the outcome
by all the comparisons. But under the hood, the math is exactly the
same. To see this, lets just quickly call the coefficients for the
perfectionism.ANOVA model.

\begin{Shaded}
\begin{Highlighting}[]
\FunctionTok{coefficients}\NormalTok{(perfectionism.ANOVA)}
\end{Highlighting}
\end{Shaded}

\begin{verbatim}
## (Intercept)   CountryUK  CountryUSA 
##  4.57699771 -0.06365015  0.08215352
\end{verbatim}

Those coefficients look familiar? They are exactly the same as we
calculated for the linear model (apart from the fact that Canada is the
reference group here, not the UK).

We will find out more about the omnibus F-test next semester when we
look at factorial ANOVA, but for now just know that it is equivalent to
the overall model variance explained by all of the levels of the
explanatory variable in the supernova output above.

\subsection{Post-hoc testing with Bonferroni
adjustment}\label{post-hoc-testing-with-bonferroni-adjustment}

Of course, the drawback of aov() is that while the omnibus test tells
that there is a difference between the countries, it doesn't tell of
where those differences are. So we must use the perfectionism.ANOVA
model to generate specific comparisons for of the multiple groups.

We have 3 comparisons in this model we need to test (i.e., UK vs USA,
USA vs Canada, and Canada vs USA). The first idea that might come to
mind is to just test each comparison separately with an independent
t-test, using an significance level of .05. At first glance, this
doesn't seem like a bad idea. However, consider a case where you don't
have 3 comparisons but you have 20 comparisons to test, and a
significance level of 0.05. What's the probability of observing at least
one significant result in those 20 tests just due to chance? High!
Specifically, with 20 comparisons being made, we have a 64\% chance of
observing at least one significant difference, even if all of the
differences are actually zero. In psychology, it's not unusual for the
number of simultaneous comparisons to be quite a bit larger than
3\ldots{} and the probability of getting a significant result simply due
to chance keeps going up as we add more comparisons.

All this is to say that running multiple comparisons, or independent
t-tests, is analogous to rolling a dice six times instead of once. The
more simultaneous tests we run the more likely we are to find a
difference even though none exists. So.. We need to adjust our thinking
and our confidence to account for the fact that we are making multiple
comparisons (i.e., simultaneous comparisons). Our significance level
must be made more conservative to account for the fact we are making
multiple simultaneous comparisons.

For this reason, methods for dealing with multiple comparisons
frequently call for adjusting the significance level in some way, so
that the probability of observing at least one significant result due to
chance remains below your desired significance level.

The Bonferroni correction does just this, and sets the significance
cut-off at α/n.~For instance, in the example above, with 20 comparisons
and significance level of = 0.05, you'd only reject a null hypothesis if
the p-value is less than 0.0025 (i.e., .05/20). In our example, we have
3 comparisons, so our p value cut-off will be .05/3 = .017. Accordingly,
for each comparison, we are looking for a t ratio that is large enough
to be considered statistically significant at the .017 level. In other
words, if we ran the three groups comparsions as independent t-tests
(like above) we would need to adjust the critical value of p to .017.

Luckly, the pairwise.t.test() function in R makes this correction for us
on each of the p vlaues for each of the comparisons. Lets go ahead and
request the pairwise comparisons from the ANOVA model of perfectionism
using this function and use the Bonferroni ajustment to the p value to
adjudicate statistical significance.

\begin{Shaded}
\begin{Highlighting}[]
\FunctionTok{pairwise.t.test}\NormalTok{(perfectionism}\SpecialCharTok{$}\NormalTok{SOP, perfectionism}\SpecialCharTok{$}\NormalTok{Country, }\AttributeTok{p.adj=}\StringTok{"bonf"}\NormalTok{) }\CommentTok{\# this function takes the outcome variable first, the explanatory variable next, and add the p value adjustment for bonferroni (i.e., p.adj = bonf).}
\end{Highlighting}
\end{Shaded}

\begin{verbatim}
## 
##  Pairwise comparisons using t tests with pooled SD 
## 
## data:  perfectionism$SOP and perfectionism$Country 
## 
##     CAN   UK   
## UK  0.664 -    
## USA 0.169 0.032
## 
## P value adjustment method: bonferroni
\end{verbatim}

The code outputs the bonferroni adjusted p values, which are more
conservative than simple t-test comparisons of 2 group means (because we
are testing 3 group means). Using the Bonferroni adjustment, only the
UK-USA comparison is statistically significant. This suggests that SOP
is higher in the US than the UK, but that there is insufficient
statistical support to distinguish between the the UK and Canada and the
US and Canada. So good news here is that even with our more conservative
bonferroni adjustment we still have empirical support for this
difference.

\subsection{Post-hoc testing with Tukey's HSD (more on this method next
term)}\label{post-hoc-testing-with-tukeys-hsd-more-on-this-method-next-term}

Another tool for post-hoc comparisons, which is not as conservative as
the Bonferroni adjustment, is Tukey's ``Honestly Significant
Difference'' (or Tukey's HSD for short). It constructs a critical HSD
value from the studentized range that the paired means must differ by to
be considered honestly significantly different (see here:
\url{http://www.real-statistics.com/statistics-tables/studentized-range-q-table/}).
To all intents and purposes, it is interpreted in the same way we have
been interpreting p values to date but it has been adjusted to account
for the increased false positive rate that is present when making
multiple comparisons. Lets go ahead and request Tukey's HSD using the
TukeyHSD() function from the perfectionism.ANOVA object.

\begin{Shaded}
\begin{Highlighting}[]
\FunctionTok{TukeyHSD}\NormalTok{(perfectionism.ANOVA)}\CommentTok{\# Now we have to use the perfectionism.ANOVA model to find out the pair of countries which differ. For this you may use the Tukey’s HSD test. Don;t worry too much about what this means for now, the output will be the same as the lm() and that is the key learning outcome.}
\end{Highlighting}
\end{Shaded}

\begin{verbatim}
##   Tukey multiple comparisons of means
##     95% family-wise confidence level
## 
## Fit: aov(formula = SOP ~ Country, data = perfectionism)
## 
## $Country
##                diff         lwr        upr     p adj
## UK-CAN  -0.06365015 -0.18644836 0.05914807 0.4386157
## USA-CAN  0.08215352 -0.01891969 0.18322673 0.1352622
## USA-UK   0.14580366  0.01226257 0.27934476 0.0287205
\end{verbatim}

Using the Tukey's HSD adjustment, only the UK-USA comparison is
statistically significant (i.e., p \textless{} .05). This suggests that
SOP is higher in the US than the UK, but that there is insufficient
statistical support to distinguish between the the UK and Canada and the
US and Canada. Exactly the same as the Bonferroni adjustment, but note
that the p-values are slightly less conservative (i.e., lower).

\subsection{Visulaisation}\label{visulaisation}

If you are interested in visualizing results of ANOVA and post-hoc tests
on the same plot (directly on the boxplots), here is a piece of code
which may be of interest to you:

\begin{Shaded}
\begin{Highlighting}[]
\DocumentationTok{\#\# set up of elements}
\NormalTok{x }\OtherTok{\textless{}{-}} \FunctionTok{which}\NormalTok{(}\FunctionTok{names}\NormalTok{(perfectionism) }\SpecialCharTok{==} \StringTok{"Country"}\NormalTok{) }\CommentTok{\# name of independent variable}
\NormalTok{y }\OtherTok{\textless{}{-}} \FunctionTok{which}\NormalTok{(}
  \FunctionTok{names}\NormalTok{(perfectionism) }\SpecialCharTok{==} \StringTok{"SOP"} \CommentTok{\# name of outcome variable}
\NormalTok{)}
\NormalTok{method1 }\OtherTok{\textless{}{-}} \StringTok{"anova"} \CommentTok{\# overall test of main effect = anova }
\NormalTok{method2 }\OtherTok{\textless{}{-}} \StringTok{"t.test"} \CommentTok{\# test of specific comparisons = t test}
\NormalTok{my\_comparisons }\OtherTok{\textless{}{-}} \FunctionTok{list}\NormalTok{(}\FunctionTok{c}\NormalTok{(}\StringTok{"CAN"}\NormalTok{, }\StringTok{"USA"}\NormalTok{), }\FunctionTok{c}\NormalTok{(}\StringTok{"USA"}\NormalTok{, }\StringTok{"UK"}\NormalTok{), }\FunctionTok{c}\NormalTok{(}\StringTok{"UK"}\NormalTok{, }\StringTok{"CAN"}\NormalTok{)) }\CommentTok{\# comparisons for post{-}hoc tests}
\DocumentationTok{\#\# building the plot}
\FunctionTok{library}\NormalTok{(ggpubr)}
\ControlFlowTok{for}\NormalTok{ (i }\ControlFlowTok{in}\NormalTok{ y) \{}
  \ControlFlowTok{for}\NormalTok{ (j }\ControlFlowTok{in}\NormalTok{ x) \{}
\NormalTok{    p }\OtherTok{\textless{}{-}} \FunctionTok{ggboxplot}\NormalTok{(perfectionism,}
      \AttributeTok{x =} \FunctionTok{colnames}\NormalTok{(perfectionism[j]), }\AttributeTok{y =} \FunctionTok{colnames}\NormalTok{(perfectionism[i]),}
      \AttributeTok{color =} \FunctionTok{colnames}\NormalTok{(perfectionism[j]),}
      \AttributeTok{legend =} \StringTok{"none"}\NormalTok{,}
      \AttributeTok{palette =} \StringTok{"npg"}\NormalTok{,}
      \AttributeTok{add =} \StringTok{"jitter"}
\NormalTok{    )}
    \FunctionTok{print}\NormalTok{(}
\NormalTok{      p }\SpecialCharTok{+} \FunctionTok{stat\_compare\_means}\NormalTok{(}\FunctionTok{aes}\NormalTok{(}\AttributeTok{label =} \FunctionTok{paste0}\NormalTok{(..method.., }\StringTok{", p{-}value = "}\NormalTok{, ..p.format..)),}
        \AttributeTok{method =}\NormalTok{ method1, }\AttributeTok{label.y =} \FunctionTok{max}\NormalTok{(perfectionism[, i], }\AttributeTok{na.rm =} \ConstantTok{TRUE}\NormalTok{)}
\NormalTok{      )}
      \SpecialCharTok{+} \FunctionTok{stat\_compare\_means}\NormalTok{(}\AttributeTok{comparisons =}\NormalTok{ my\_comparisons, }\AttributeTok{method =}\NormalTok{ method2, }\AttributeTok{label =} \StringTok{"p.format"}\NormalTok{) }
\NormalTok{    )}
\NormalTok{  \}}
\NormalTok{\}}
\end{Highlighting}
\end{Shaded}

\begin{verbatim}
## Warning: The dot-dot notation (`..method..`) was deprecated in ggplot2 3.4.0.
## i Please use `after_stat(method)` instead.
## i The deprecated feature was likely used in the ggpubr package.
##   Please report the issue at <https://github.com/kassambara/ggpubr/issues>.
## This warning is displayed once every 8 hours.
## Call `lifecycle::last_lifecycle_warnings()` to see where this warning was
## generated.
\end{verbatim}

\includegraphics{Notebook_empty_files/figure-latex/unnamed-chunk-30-1.pdf}

\subsection{Effect size}\label{effect-size-1}

We can also calculate effect sizes for the group difference just like we
did above.

\begin{Shaded}
\begin{Highlighting}[]
\NormalTok{perfectionism }\SpecialCharTok{\%\textgreater{}\%}
\FunctionTok{cohens\_d}\NormalTok{(SOP }\SpecialCharTok{\textasciitilde{}}\NormalTok{ Country)}
\end{Highlighting}
\end{Shaded}

\begin{verbatim}
## # A tibble: 3 x 7
##   .y.   group1 group2 effsize    n1    n2 magnitude
## * <chr> <chr>  <chr>    <dbl> <int> <int> <ord>    
## 1 SOP   CAN    UK       0.224    73    24 small    
## 2 SOP   CAN    USA     -0.463    73    42 small    
## 3 SOP   UK     USA     -0.537    24    42 moderate
\end{verbatim}

According to Cohen (1988, 1992), the effect size is low if the value of
d varies around 0.2, medium if d varies around 0.5, and large if d
varies more than 0.8. The UK-USA difference is a large effect and other
two comparisons are small effects.

\subsection{How its written}\label{how-its-written-2}

A one-way between-subjects ANOVA was conducted to compare the effect of
country on levels of perfectionism. There was a significant effect of
country on perfectionism at the p\textless.05 level for the three
conditions {[}F(2, 136) = 3.65, p = 0.03{]}. Post hoc comparisons using
the Tukey HSD test indicated that the mean perfectionism score for the
US (M = 4.66, SD = 0.16) was significantly larger than the mean
perfectionism score for the UK (M = 4.51, SD = 0.35, d=.54). However,
the mean perfectionism score for Canada (M = 4.58, SD = 0.20) did not
significantly differ from the mean perfectionism scores for the US
(d=.46) and the UK (d = .22).

\subsection{Exercise}\label{exercise}

Okay, your turn to run some of the tests that we have described here. To
do that, let's once again return to our WVS data and specifically the
comparison of perceptions of choice between Australia and Japan that we
made last week. Then, we compared the distributions using z-scores.
Let's now go one step further and test the differences in perceptions of
choice using a two-parameter linear model or independent t-test.

Before we start, lets just clarify the research question for our
two-parameter linear model and null hypothesis we are testing:

\emph{Research Question} - Do perceptions of choice differ between
people from Australia and Japan?

\emph{Null Hypothesis} - The difference between perceptions of choice in
Australia and Japan will be zero.

\subsection{Task 1 - Loading the
dataframe}\label{task-1---loading-the-dataframe}

Load the WVS dataframe. The WVS data is avalaible in the MT9 workshop
folder. Load it into the R environment like you have been doing in
previous weeks (i.e., click on it, change name to WVS and copy+paste the
code below).

\begin{Shaded}
\begin{Highlighting}[]
\NormalTok{WVS }\OtherTok{\textless{}{-}} \FunctionTok{readRDS}\NormalTok{(}\StringTok{"C:/Users/currant/Dropbox/Work/LSE/PB130/MT9/Workshop/WV6\_Data\_R\_v20180912.rds"}\NormalTok{)}

\NormalTok{WVS}\SpecialCharTok{$}\NormalTok{V2A }\OtherTok{\textless{}{-}} \FunctionTok{as\_character}\NormalTok{(WVS}\SpecialCharTok{$}\NormalTok{V2A) }\CommentTok{\# maintain this code for the country variable}
\end{Highlighting}
\end{Shaded}

\subsection{Task 2 - Selecting the
variables}\label{task-2---selecting-the-variables}

I looked in the codebook for WVS and the perceptions of choice variable
is listed as V55. So lets look for the variable and select it, then head
the new dataframe entitled WVSchoice. Lets also select the country code
variable, V2A, so we can use this to group the choice scores later. I
have left some of the below code blank for you to fill the gap.

\begin{Shaded}
\begin{Highlighting}[]
\NormalTok{WVSchoice }\OtherTok{\textless{}{-}} 
\NormalTok{  WVS }\SpecialCharTok{\%\textgreater{}\%} 
  \FunctionTok{select}\NormalTok{(V2A, V55) }\CommentTok{\# SELECT THE VARIABLES OF INTEREST HERE}
\FunctionTok{head}\NormalTok{(WVSchoice)}
\end{Highlighting}
\end{Shaded}

\subsection{Task 3 - Filter the
variables}\label{task-3---filter-the-variables}

Just like other variables in the WVS dataset, we need to edit the
dataframe since the choice variable has values we can't interpret like
-5. While we are at it, lets also filter out Australia and Japan as the
countries of interest. I have left some of the below code blank for you
to fill the gap.

\begin{Shaded}
\begin{Highlighting}[]
\NormalTok{WVSchoice }\OtherTok{\textless{}{-}} 
\NormalTok{WVSchoice }\SpecialCharTok{\%\textgreater{}\%} 
\FunctionTok{filter}\NormalTok{(V55 }\SpecialCharTok{\textgreater{}=} \StringTok{"1"} \SpecialCharTok{\&}\NormalTok{ V2A }\SpecialCharTok{==} \FunctionTok{c}\NormalTok{(}\StringTok{"Australia"}\NormalTok{, }\StringTok{"Japan"}\NormalTok{)) }\CommentTok{\# Fill in the question marks }
\NormalTok{WVSchoice}
\end{Highlighting}
\end{Shaded}

\subsection{Task 4 - Examine the
descriptives}\label{task-4---examine-the-descriptives}

Use the favstats function to inspect the means and SD of the choice
scores for Australia and Japan.

\begin{Shaded}
\begin{Highlighting}[]
\FunctionTok{favstats}\NormalTok{(}\SpecialCharTok{\textasciitilde{}}\NormalTok{V55, V2A, }\AttributeTok{data=}\NormalTok{WVSchoice) }\CommentTok{\# Fill in the question mark}
\end{Highlighting}
\end{Shaded}

What is the mean score for Australia?

What is the mean score for Japan?

Calculate the mean difference

\begin{Shaded}
\begin{Highlighting}[]
\FloatTok{7.742779}    \SpecialCharTok{{-}} \FloatTok{5.811899}  
\end{Highlighting}
\end{Shaded}

What is the mean difference?

\subsection{Task 5 - Setting up the linear
model}\label{task-5---setting-up-the-linear-model}

Now lets set our up our linear model to examine the statistical
significance of that mean difference. We are using the lm() function to
test a two-parameter model because we now have two means to estimate -
one for Australia and one for Japan. As I mentioned above, R actually
does the dummy coding for us in the lm() function so we don't need to
dummy code the country variable. For noting, R will code Australia as 0
and Japan as 1 (just because thats the order they appear in the data
set). Lets just build the model with V2A as the explanatory factor.

\begin{Shaded}
\begin{Highlighting}[]
\NormalTok{choice.model }\OtherTok{\textless{}{-}} \FunctionTok{lm}\NormalTok{(V55 }\SpecialCharTok{\textasciitilde{}} \DecValTok{1} \SpecialCharTok{+}\NormalTok{ V2A, }\AttributeTok{data=}\NormalTok{WVSchoice) }\CommentTok{\# fill in the variable needed for question mark}
\FunctionTok{summary}\NormalTok{(choice.model)}
\end{Highlighting}
\end{Shaded}

What is the mean score for Australia (hint: intercept)?

What is the incremenent needed to find the Japan mean or - put
differently - the mean difference?

What is the t-ratio?

What is the p value(remember that e-values mean zero)?

Do we accept or reject the null hypothesis?

Why?

\subsection{Task 6 - Paritioning
variance}\label{task-6---paritioning-variance}

We've established whether there is a statistically significant mean
difference, but how much variance in choice is explained by country? We
can answer this question using the supernova() function in R, which
breaks down the variance in choice due to the model (ie., country) and
error (i.e., variance left over once we've subtracted out the model)

\begin{Shaded}
\begin{Highlighting}[]
\FunctionTok{supernova}\NormalTok{(choice.model) }\CommentTok{\# add the choice model you want to examine the variance}
\end{Highlighting}
\end{Shaded}

What is the total variance or sum of squares for perceptions of choice?

How is that sum of squares partitioned between the model?

And the Error?

What is the percentage of variance explained in by the model?

What is the F ratio?

What is the p value?

What do we conclude in relation to our research question?

\end{document}
